problem:  
Let \(\{(X, \omega_t, \Omega_t)\}_{t \in [0,1]}\) be a smooth family of Calabi–Yau threefolds on a fixed differentiable manifold \(X\), such that the variation of Hodge structure on \(H^3(X,\mathbb{C})\) has monodromy operator \(T\) when the family is continued along a closed loop in the complex moduli.  
Suppose \(L_0 \subset (X,\omega_0,\Omega_0)\) is a compact embedded special Lagrangian submanifold representing a class \(\alpha_0 \in H_3(X,\mathbb{Z})\).  

Let \(\mathcal{M}(t)\) denote the local McLean moduli space of special Lagrangians near the transported submanifold \(L_t\) (viewed under a smooth identification of the underlying differentiable manifold).  Assume that for each \(t\in [0,1]\), there exists an isolated special Lagrangian representative \(L_t\) of phase \(\theta_t = \arg \int_{\alpha_t}\Omega_t\), where \(\alpha_t\) is the homology class obtained from \(\alpha_0\) by parallel transport under the Gauss–Manin connection.

**Problem:**  
When the family \((X,\omega_t,\Omega_t)\) is analytically continued around the loop and returned to \(t=1\) with \((\omega_1,\Omega_1)=(\omega_0,\Omega_0)\), the monodromy operator acts as \(T:\alpha_0 \mapsto T\alpha_0\).  

Is it possible—possibly after enlarging the moduli space to include calibrated currents or varifold limits—to define a canonical “monodromy transform”  
\[
\mathfrak{T} : \overline{\mathcal{M}}_{\mathrm{SLag}}(\alpha_0) \longrightarrow \overline{\mathcal{M}}_{\mathrm{SLag}}(T\alpha_0)
\]
such that:  
1. on smooth points representing non-singular special Lagrangians \(L_0\), the image \(\mathfrak{T}(L_0)\) corresponds (in an appropriate weak sense) to the limit of analytic continuation of the family \(L_t\) along the loop; and  
2. this correspondence preserves calibrated volume, up to the phase rotation induced by the monodromy of \(\int_{\alpha_t}\Omega_t\)?

Equivalently, can one construct a geometric realization of period monodromy—classically acting on \(H_3(X,\mathbb{Z})\)—as an *analytic transport* acting on the compactified special Lagrangian moduli space, potentially exchanging homology classes of distinct phases through degeneration limits and reemergence of special Lagrangians?

Why is it a "good" problem:  
This problem makes rigorous the conceptual question of continuity and monodromy of the special Lagrangian condition under loops in the complex moduli of Calabi–Yau manifolds.  It preserves mathematical precision by focusing on the tangible goal of constructing a well-defined monodromy transform on the (compactified) moduli of calibrated submanifolds or currents.

Its significance lies in establishing a direct link between three foundational structures:
- **Geometric analysis:** the nonlinear elliptic PDEs governing calibrated geometry and their compactness/limit behavior.
- **Hodge theory:** the Gauss–Manin connection and period monodromy.
- **Mirror symmetry:** realization of wall-crossing and stability transitions as real-analytic geometric transformations.

A positive resolution would provide the first rigorous analytic manifestation of period-monodromy action on the *real side* of mirror symmetry, showing how topological and analytic data of Calabi–Yau geometry interact through special Lagrangians.  
The problem is straightforward to state—a search for a canonical transformation induced by loop monodromy—but demands deep advances in calibrated geometry, moduli compactification, and analytic continuation of nonlinear PDE solutions, making it an elegant and truly “good” research problem.